\documentclass[12pt]{article}
\usepackage[T1]{fontenc}
\usepackage[utf8]{inputenc} 
\usepackage[a4paper]{geometry}
\usepackage{gensymb}
\geometry{a4paper, margin=1in}
\usepackage{natbib}
\usepackage{hyperref}
\usepackage{booktabs}
\usepackage{threeparttable}
\usepackage{amssymb}
\usepackage{amsmath}
\title{Assignment 2}
\author{Dmitrii Kuptsov}
% \date{March 2025}
\begin{document}
\maketitle
\section*{Task 1}

\subsection*{Task 1.a}
\[
d_e = 24, \quad d_m = 3.14606742
\]

\subsection*{Task 1.b}

\subsubsection*{Task 1.b.i}

If every observation is multiplied by 10, average will be also multiplied by 10. However, the average of $v_1$ is 0, so nothing changes in $\bar{v}$:

\[
\bar{v} =
\begin{bmatrix}
0 \\[4pt]
1 \\[4pt]
1
\end{bmatrix}
\]


$v_0$ is a single observation, so its first component will be multiplied by 10, nothing else changes:


\[
\bar{v_0} =
\begin{bmatrix}
20 \\[4pt]
5 \\[4pt]
3
\end{bmatrix}
\]

As for S, it will look like this:

\[
S =
\begin{bmatrix}
500 & 10 & 20 \\[4pt]
10 & 6 & 1 \\[4pt]
20 & 1 & 7
\end{bmatrix}.
\]


$S_{11}$ element is a variance of $v_1$. $v_1$ is increased in 10 times, so the variance is increased in $10^2 = 100$ times. $S_{12} = S_{21}$ and  $S_{13} = S_{31}$ are the covariances of $v_1$ with $v_2$ and $v_3$ correspondingly. $v_2$ and $v_3$ do not change, so the increase of $v_1$ in 10 times means that covariances increase in 10 times.

Other elements of S does not include $v_1$, so they are unchanged.

\subsubsection*{Task 1.b.ii}

\[
d_e = 420, \quad d_m = 3.14606742
\]

$d_e$ increased, which is reasonable, because it is simply a sum of squared distances for each component.

$d_m$ remains unchanged, because it is additionally normalized using the variances and covariances of the components.

\section*{Task 2}

\subsection*{Task 2.a}

First, we can represent X as the vector of column-vectors:

\[X = [1, ln(X_2), X_3]\]

Similarly with $X^*$:

\[X^* = [1, ln(a\cdot X_2), b\cdot X_3]\]

If X and $X^*$ are $n \times 3$, A should be $3\times 3$. Moreover, to get $i^{th}$ element of the first column of $X^*$ we need to multiply the $i^{th}$ row of X with the first column of A. Therefore, in the first column of A the first element should be 1, and other elements should be zeros.

Similarly, the third element of the third column should be b, and other elements in this column must be zeros.

As for the second column, it should transform from $ln(X_2)$ to $ln(a\cdot X_2) = a + ln(X_2)$, so in the first row it will be a (adding constant), in the second it will be 1, and third is 0.

Therefore:

\[
A =
\begin{bmatrix}
1 & a & 0 \\[4pt]
0 & 1 & 0 \\[4pt]
0 & 0 & b
\end{bmatrix}.
\]

\subsection*{Task 2.b}

\[
\begin{aligned}
\hat{\beta}^*
&= (X^{*\top} X^*)^{-1} X^{*\top} y \\[4pt]
&= \big((X A)^{\top} X A\big)^{-1} (X A)^{\top} y \\[4pt]
&= (A^{\top} X^{\top} X A)^{-1} A^{\top} X^{\top} y \\[4pt]
&= A^{-1} (X^{\top} X)^{-1} (A^{\top})^{-1} A^{\top} X^{\top} y \\[4pt]
&= A^{-1} (X^{\top} X)^{-1} X^{\top} y
\\[4pt]
&= A^{-1} \hat{\beta}.
\end{aligned}
\]

So, the relation of $\hat{\beta}^*$ and $\hat{\beta}$ equals to $A^{-1}$.

\[
A^{-1} =
\begin{bmatrix}
1 & -a & 0 \\[4pt]
0 & 1 & 0 \\[4pt]
0 & 0 & \dfrac{1}{b}
\end{bmatrix}
\]

Therefore, $\hat{\beta_1^*} = \hat{\beta_1} - a\cdot \hat{\beta_2}$. This seems reasonable, because we need to adjust the constant (the first parameter) for the change in second parameter. The multiplication of the second parameter increased the total sum of RHS by a, so we need to subtract it from the constant term to balance the sides.

The second parameter remains unchanged, because the change in it was constant, and we have accounted for that in the previous step.

The third element is multiplied by $\frac{1}{b}$. Every value corresponding to this parameter was multiplied by b, and it is included linearly in the regression, so it is only reasonable that we adjust the coefficient by the inverse value.


\section*{Task 3}

\subsection*{Task 3.a}

\[
\beta_1 = 4.9952, \quad \beta_2 = 0.0709
\]


\subsection*{Task 3.b}

\[
\beta_1 = 4.9952, \quad \beta_2 = 0.0709
\]



\end{document}